% Base Font Size
%   All other font sizes (for headings etc.) are adjusted accordingly

\KOMAoptions{fontsize=12pt}

% Depth of table of contents
%   0 displays only parts and chapters, 1 displays parts, chapters and sections, etc.

\setcounter{tocdepth}{1}

% Citation style
%   for a reference see http://mirrors.ctan.org/macros/latex/contrib/biblatex/doc/biblatex.pdf#subsection.3.3

\usepackage[style=numeric]{biblatex}

% Allow long URLs in the bibliography to break across lines
\setcounter{biburlnumpenalty}{9000}
\setcounter{biburllcpenalty}{9000}
\setcounter{biburlucpenalty}{9000}

% Drop eprint fields that are raw URLs (ACL Anthology dumps put full PDF
% links there, which print as unbreakable text); real arXiv IDs are kept
\DeclareSourcemap{
  \maps[datatype=bibtex]{
    \map{
      \step[fieldsource=eprint, match=\regexp{^https?://}, final]
      \step[fieldset=eprint, null]
    }
    % ISBNs are not printed by the numeric style; proceedings dumps carry
    % malformed ones that trigger datamodel warnings
    \map{
      \step[fieldset=isbn, null]
    }
    % Drop legacy month ranges like "13--19 Jul" that cannot sort
    \map{
      \step[fieldsource=month, match=\regexp{[\s]}, final]
      \step[fieldset=month, null]
    }
  }
}

% Print the whole bibliography?
%   By default, LaTeX prints just those bibliography items that have been cited.
%   Uncomment the following line to print the whole bibliography.
%
%   \nocite{*}


% --- Encoding and fonts ---------------------------------------------
\usepackage[utf8]{inputenc}
\usepackage[T1]{fontenc}
\usepackage{microtype}

% --- Page layout and graphics ---------------------------------------
\usepackage{geometry}
\usepackage{graphicx}
\usepackage{wrapfig}
\usepackage{xcolor}

% --- Maths ----------------------------------------------------------
\usepackage{amsfonts}
\usepackage{nicefrac}

% --- Tables and floats ----------------------------------------------
\usepackage{booktabs}
\usepackage{tabularx}
\usepackage{array}
\usepackage{multirow}
\usepackage{float}
\usepackage{placeins}
\usepackage{rotating}

% --- Lists and boxes ------------------------------------------------
\usepackage{enumitem}
\usepackage{framed}
\usepackage[most]{tcolorbox}

% --- Verbatim with automatic line breaking ---------------------------
\usepackage{fvextra}


% --- Colours for the case-study boxes -------------------------------
\definecolor{HighVulnFrame}{RGB}{200, 80, 60}
\definecolor{HighVulnBack}{RGB}{255, 240, 235}

\definecolor{LLMFrame}{RGB}{60, 100, 180}
\definecolor{LLMBack}{RGB}{240, 245, 255}

\definecolor{PromptFrame}{RGB}{80, 100, 100}
\definecolor{PromptBack}{RGB}{245, 245, 245}

\definecolor{promptframe}{gray}{0.55}


% --- Case-study boxes ------------------------------------------------
\newtcolorbox{profilebox}[1][]{
  enhanced,
  title=\textbf{User Profile (High Vulnerability)},
  colback=HighVulnBack,
  colframe=HighVulnFrame,
  coltitle=white,
  fonttitle=\bfseries\sffamily,
  boxrule=1.2pt,
  arc=3mm,
  #1
}

\newtcolorbox{promptbox}[1][]{
  enhanced,
  title=\textbf{User Prompt},
  colback=PromptBack,
  colframe=PromptFrame,
  coltitle=white,
  fonttitle=\bfseries\sffamily,
  fontupper=\ttfamily\small,
  boxrule=1.2pt,
  arc=3mm,
  #1
}

\newtcolorbox{responsebox}[1][]{
  enhanced,
  breakable,
  title=\textbf{LLM Response},
  colback=LLMBack,
  colframe=LLMFrame,
  coltitle=white,
  fonttitle=\bfseries\sffamily,
  fontupper=\small,
  boxrule=1.2pt,
  arc=3mm,
  #1
}


% --- Appendix prompt boxes -------------------------------------------
\DefineVerbatimEnvironment{promptcode}{Verbatim}{%
  fontsize=\small,
  breaklines=true,
  breakanywhere=true,
  breakindent=1.5em,
  breaksymbolleft=\textcolor{promptframe}{\small\ensuremath{\hookrightarrow}},
  breaksymbolright={},
  frame=single,
  framerule=0.4pt,
  rulecolor=\color{promptframe},
  framesep=8pt,
  xleftmargin=0pt,
  xrightmargin=0pt
}

\fvset{breaklines=true,breakanywhere=true}

\newcommand{\promptmeta}[4]{%
  \begingroup\small
  \begin{tabularx}{\linewidth}{@{}>{\bfseries}l X@{}}
    \toprule
    Purpose  & #1 \\
    Model    & #2 \\
    Settings & #3 \\
    Output   & #4 \\
    \bottomrule
  \end{tabularx}
  \endgroup
  \vspace{1ex}
}

\newcommand{\widetable}[1]{%
  \begingroup
  \footnotesize
  \setlength{\tabcolsep}{3.5pt}
  \renewcommand{\arraystretch}{1.05}
  \input{#1}%
  \endgroup
}

\providecommand{\subcell}[1]{\begin{tabular}[t]{@{}l@{}}#1\end{tabular}}
\newcolumntype{P}[1]{>{\raggedright\arraybackslash}p{#1}}


% --- URLs and hyperlinks (load last) ----------------------------------
\usepackage{url}
\Urlmuskip=0mu plus 1mu
\def\UrlBreaks{\do\/\do-\do_\do.}
\usepackage{hyperref}
\hypersetup{hidelinks} % no link colours or border boxes (print-friendly); links stay clickable
\usepackage{xurl}